\subsection{Characterizations of Dedekind domains}

THEOREM 1. Let A be an integral domain and K its field of fractions. The following conditions are equivalent:

\begin{enumerate}
\def\labelenumi{(\alph{enumi})}
\item
  A is a Dedekind domain;
\item
  A is a Krull domain and every non-improper valuation on K which is positive on A is equivalent to an essential valuation of A;
\item
  A is a Krull domain and every fractional ideal \(\Im \neq (0)\) of A is divisorial;
\item
  every fractional ideal \(\Im \neq (0)\) of A is invertible;
\item
  A is a Noetherian integrally closed domain and every non-zero prime ideal of A is maximal;
\item
  A is Noetherian and, for every maximal ideal m of A, \(A_{m}\) is either a field or a discrete valuation ring;
\item
  A is Noetherian and, for every maximal ideal m of A, \(A_{m}\) is a principal ideal domain.
\end{enumerate}

We show first the equivalence of (a) and (b). Corollary 2 to Theorem 3, §1, no. 6, shows immediately that (a) implies (b). Conversely, (b) implies (a), since, for every prime ideal p of A, there exists a valuation ring of K which dominates \(A_{p}\) (Chapter VI, § 1, no. 2, Corollary to Theorem 2).

The remainder of the proof is carried out by proving the following implications:

\[
(a) \Rightarrow (c) \Rightarrow (d) \Rightarrow (e) \Rightarrow (f) \Rightarrow (g) \Rightarrow (a).
\]

If A is a Dedekind domain and b is a non-zero fractional ideal, then \(bA_{p} = \tilde{b}A_{p}\) for every maximal ideal p (§ 1, no. 4, Proposition 7) and hence \(b = \tilde{b}\) (Chapter II, § 3, no. 3, Corollary 3 to Theorem 1); thus (a) implies (c). We now show that (c) implies (d). If (c) holds, the mapping \(a \mapsto div a\) is a bijection of I(A) onto D(A) (cf.~§ 1, no. 1); as it is a homomorphism (§ 1, no. 2) and D(A) is a group, every element of I(A) is invertible.

We show that (d) implies (e). If (d) holds, every integral ideal \(\neq(0)\) of A is finitely generated (Chapter II, § 5, no. 6, Theorem 4) and hence A is Noetherian; as I(A) is a group, D(A) is a group and A is therefore completely integrally closed (§ 1, no. 2, Theorem 1). Finally, if p is a non-zero prime ideal of A and m is a maximal ideal of A containing p, the ring \(A_{m}\) is a principal ideal domain (Chapter II, § 5, no. 6, Theorem 4); as \(pA_{m}\) is prime and non-zero, necessarily \(pA_{m}=mA_{m}\) (a principal ideal domain being a Dedekind domain) whence p=m (Chapter II, § 2, no. 5, Proposition 11) and p is maximal.

We now show that (e) implies (f). If m is a maximal ideal of A and (e) holds, \(A_{m}\) is an integrally closed Noetherian domain and its maximal ideal \(mA_{m}\) is, either (0), or the only non-zero prime ideal of \(A_{m}\) ; hence \(A_{m}\) is a field or a discrete valuation ring by Proposition 11 of §1, no. 7.

The fact that (f) implies (g) is obvious.

We show finally that (g) implies (a). As A is the intersection of the \(A_{m}\) , where m runs through the set of maximal ideals (Chapter II, §3, no. 3, Corollary 4 to Theorem 1), (g) implies that A is integrally closed and Noetherian and hence that A is a Krull domain (§1, no. 3, Corollary to Theorem 2). On the other hand, it can be shown that every non-zero prime ideal of A is maximal as in the proof that (d) \(\Rightarrow\) (e).

PROPOSITION 1. A semi-local Dedekind domain is a principal ideal domain.

Let A be a semi-local Dedekind domain, K its field of fractions, \(p_{1}, \ldots, p_{n}\) its maximal ideals and \(v_{1}, \ldots, v_{n}\) the corresponding essential valuations; these are the only essential valuations of A. Let a be a non-zero integral ideal of A. Since it is divisorial, there exists (§ 1, no. 4, Proposition 5) integers \(q_{1}, \ldots, q_{n}\) such that a is the set of \(x \in K\) such that \(v_{i}(x) \geqslant q_{i}\) for \(1 \leqslant i \leqslant n\) . Let \(x_{0}\) be an element of K such that \(v_{i}(x_{0}) = q_{i}\) for \(1 \leqslant i \leqslant n\) (Chapter VI, § 7, no. 2, Corollary 1 to Theorem 1). Then a is the set of \(x \in K\) such that \(v_{i}(xx_{0}^{-1}) \geqslant 0\) for \(1 \leqslant i \leqslant n\) . Thus \(a = Ax_{0}\) .

If A is a Dedekind domain, it has been seen, in the proof of Theorem 1, that the group D(A) of divisors of A is identified with the group I(A) of fractional ideals \(a \neq (0)\) (as A is Noetherian, every non-zero fractional ideal is finitely generated). The divisor class group C(A) of A ( \(\S 1\) , no. 2) is then identified with the group of classes of ideals \(\neq 0\) of A (defined in Chapter II, \(\S 5\) , no. 7).

